%% LyX 2.5.1 created this file.  For more info, see https://www.lyx.org/.
%% Do not edit unless you really know what you are doing.
\documentclass[brazil]{article}
\usepackage[T1]{fontenc}
\usepackage[utf8]{luainputenc}

\makeatletter

%%%%%%%%%%%%%%%%%%%%%%%%%%%%%% LyX specific LaTeX commands.

\makeatother

\usepackage{babel}
\begin{document}

\subsection{Há desigualdade (e isso importa)}

Não há segredo sobre a existência de desigualdade no capitalismo. Muitas fontes e estudos diferentes podem ser consultadas. Para desenvolver essa discussão eu escolhi um relatório de 2019 elaborado pela ONU \cite{un2019sustainable}.

Os Objetivos de Desenvolvimento Sustentável foram concebidos no âmbito da Agenda 2030 para o Desenvolvimento Sustentável e representam uma das maiores esperanças para um futuro sustentável, mas enfrentam desafios significativos. Nas palavras do prefácio do relatório:
\begin{quote}
Níveis persistentemente elevados de desigualdade perpetuam a incerteza e a insegurança entre as pessoas, reforçando divisões e minando a confiança nas instituições e no governo.
\end{quote}
O enfraquecimento do crescimento econômico global e o aumento da desigualdade econômica representam uma séria ameaça à implementação dos Objetivos de Desenvolvimento Sustentável. Para que esses objetivos sejam alcançados, a comunidade internacional deve olhar além do mero crescimento econômico. Um aumento no PIB não reflete necessariamente uma melhoria na qualidade de vida.

Persiste um nível elevado de insegurança laboral no mundo: a maior parte da população está envolvida em trabalho informal ou vulnerável\footnote{Os trabalhadores vulneráveis  ocupam empregos pouco qualificados e mal remunerados, com pouca segurança no emprego. Muitas vezes, são explorados pelos seus empregadores e forçados a trabalhar longas horas por baixos salários ou em condições perigosas\cite{unison2026vulnerable}.}, e um em cada três adultos enfrenta riscos crônicos de desemprego.

É necessário reduzir a desigualdade em todas as suas formas e manifestações. A Agenda 2030 adotou o lema “não deixar ninguém para trás” como princípio orientador. No entanto, na prática, há poucas evidências de que o mundo esteja vencendo a luta contra a desigualdade. Essa desigualdade impede o desenvolvimento sustentável de diversas maneiras: por meio da estagnação econômica, da limitação da mobilidade social e do desenvolvimento humano, e da erosão da confiança nas instituições.

Na prática, os efeitos da desigualdade podem ser observados, por exemplo, no fato de que conflitos violentos frequentemente surgem de desigualdades entre grupos e da insegurança habitacional. Portanto, é necessário ampliar as oportunidades econômicas e garantir moradia, especialmente para os grupos mais marginalizados e vulneráveis.

Nos países em desenvolvimento, a população gasta cerca de 80\% da renda total com necessidades básicas (alimentação, moradia, saúde e energia), enquanto mesmo nos países desenvolvidos, o gasto com recursos básicos chega a quase 60\%. Uma em cada cinco pessoas no mundo vive em moradias inadequadas e um bilhão de pessoas ainda vivem em assentamentos informais. Para enfrentar essa situação, é necessário agir em todas as dimensões, fortalecendo as medidas de proteção social e promovendo a redistribuição de riqueza, renda e oportunidades.

É evidente que o mundo tolerou o aumento da desigualdade por tempo demais e agora precisa encarar suas consequências, que corroem a solidariedade e a coesão social. Os esforços nacionais e internacionais devem ser direcionados para conter e reduzir a desigualdade, e isso só pode ser alcançado por meio de medidas regulatórias. Uma coisa é certa: o crescimento econômico por si só não tirará todos da pobreza, e as políticas devem ir além do crescimento do PIB. Os melhores resultados obtidos até agora na busca pelo desenvolvimento provêm de uma combinação de políticas nacionais, como as voltadas para o mercado de trabalho, e políticas de redistribuição, que protegem os mais vulneráveis e marginalizados.

Na maioria dos países, o crescimento salarial não acompanhou o ritmo da produtividade, em grande parte devido à ausência de políticas proativas. Os salários permaneceram estagnados para os 50\% mais pobres da sociedade, enquanto aumentaram substancialmente para os 10\% e 1\% mais ricos. Ao mesmo tempo, o custo de vida subiu, afetando desproporcionalmente as camadas mais baixas da pirâmide social. Dados coletados em 2019 mostram que, nos Estados Unidos, 15\% da população vive em situação de alta insegurança financeira e 25\% em situação de insegurança moderada. O relatório também aponta que a insegurança financeira influencia o comportamento eleitoral, favorecendo a extrema direita.

Quanto ao resto do mundo, mais de dois terços dos países estão a registrar um aumento da desigualdade de rendimento e riqueza. Existe um mecanismo de reforço entre o emprego vulnerável e os níveis de desigualdade: por um lado, os elevados níveis de emprego vulnerável reforçam a desigualdade persistente, enquanto, por outro, a elevada desigualdade contribui para a expansão do trabalho vulnerável.

Diante deste cenário, o setor público não pode ser excluído do debate, uma vez que é um fato observável que a concentração de riqueza resulta também numa concentração de poder político. Nos Estados Unidos, as instituições públicas atuam frequentemente promovendo políticas que favorecem uma pequena e rica fração da população em detrimento da maioria. Esta elevada concentração de poder político permite aos ricos impor políticas que servem os seus próprios interesses, mesmo quando tais políticas não beneficiam o público em geral.

Desta forma, a concentração de poder no topo da pirâmide social permite um processo de tomada de decisões que ignora e não satisfaz as necessidades das pessoas fora deste segmento privilegiado. Buscando proteger seus próprios interesses, essa fração minoritária da população, aqueles que detêm o poder político e econômico, tem tanto o incentivo quanto os meios para bloquear ativamente a implementação de medidas que reduzam a desigualdade, criando, na prática, uma armadilha da desigualdade. Os níveis mais altos de desigualdade estão associados a menor sindicalização e menor proteção trabalhista no setor público.

Ao analisarmos os dados globais, podemos observar que a desigualdade entre países, medida pelo coeficiente de Gini, diminuiu. No entanto, esse resultado deve-se em grande parte às enormes populações dos países asiáticos, cuja riqueza e renda médias per capita historicamente estiveram abaixo da média global. À medida que países como a China e a Índia convergem para as médias globais, o coeficiente de Gini global diminui automaticamente, mesmo que a desigualdade no interior destes países aumente.

Por outro lado, a desigualdade dentro dos países, também medida pelo índice de Gini, aumentou em nações que, juntas, representam mais de dois terços da população mundial. É necessário, ainda, interpretar o índice de Gini criticamente: países que registraram uma queda no coeficiente de Gini podem experientar simultaneamente um aumento na parcela da riqueza detida pelos 10\% (ou 1\%) mais ricos, visto que o índice de Gini é menos sensível a variações nos extremos da distribuição.

O relatório destaca ainda que mesmo os países que apresentam uma redução na desigualdade, segundo diferentes índices, têm experimentado um crescimento salarial e uma criação de empregos que permanecem muito lentos para permitir que aqueles na base da pirâmide de renda escapem da pobreza.

Outra consequêncis da desigualdade inclui seu impacto restritivo sobre a mobilidade social. A elevada desigualdade está intimamente ligada à baixa mobilidade social, minando a narrativa meritocrática do capitalismo, uma vez que as famílias mais pobres enfrentam barreiras financeiras que as impedem de investir adequadamente na educação e nas oportunidades de ascensão social de seus filhos.

A desigualdade também se manifesta na forma como as crises financeiras afetam as diferentes classes sociais de maneira desigual. Enquanto a classe média sofre desproporcionalmente com as recessões econômicas, a classe alta não só tem maior capacidade de se preparar para crises, como muitas vezes recebe até mesmo compensação financeira dos bancos. Enquanto isso, a classe baixa tende a permanecer isolada geográfica e economicamente dos mercados nacionais e globais devido à sua posição estrutural dentro do sistema econômico.

A globalização do comércio tem sido frequentemente citada como um fator-chave no aumento da desigualdade tanto em países desenvolvidos quanto em desenvolvimento. No entanto, essa afirmação não é fortemente sustentada por evidências empíricas. Dados comparativos entre países nas últimas décadas mostram apenas uma correlação limitada entre a abertura comercial e a desigualdade interna. De acordo com estudos recentes, embora a liberalização comercial desempenhe um papel na desigualdade de renda, seu impacto é secundário ao do desenvolvimento tecnológico, da abertura financeira e da profundidade financeira. Os principais fatores determinantes são, na verdade, políticas públicas ineficazes, como o enfraquecimento do poder de negociação dos trabalhadores, o desenvolvimento educacional inadequado e o desmantelamento das redes de proteção social.

Um fato alarmante mostra que 40\% da população nos países da OCDE é economicamente vulnerável, o que significa que não possui capital financeiro para viver por três meses acima da linha da pobreza caso fiquem sem renda. Há muito mais pessoas economicamente vulneráveis  do que aquelas classificadas como pobres em termos unicamente de renda. Situações de insegurança estão diretamente ligadas aos níveis de desigualdade: a desigualdade muito alta intensifica os sentimentos de insegurança e incerteza. Por outro lado, o aumento da insegurança tem impactos diferentes em diferentes grupos populacionais, o que pode agravar ainda mais a desigualdade.

Por exemplo, a insegurança econômica afeta o desempenho cognitivo. Experimentos em laboratório mostram que as pessoas têm desempenhos diferentes dependendo da diferença entre o que precisam e os recursos disponíveis para atingir seus objetivos. Quanto maior essa diferença, mais a consciência da situação captura a atenção e gera pensamentos intrusivos, prejudicando a capacidade cognitiva. Aqueles que enfrentam insegurança econômica sofrem uma espécie de " imposto cognitivo" , que afeta seu bem-estar e produtividade.

A desigualdade também impacta a resiliência das famílias a eventos climáticos, que depende fortemente das condições socioeconômicas. Ela mina a confiança pública nos governos, nas instituições e até mesmo entre os indivíduos, criando um ambiente político favorável a candidatos considerados ``\textit{outsiders}'' ou ``fora do sistema'', especialmente populistas e nacionalistas.

Essa divisão entre elites e trabalhadores comuns em nível nacional também se projeta internacionalmente, com forças nacionalistas capitalizando sobre o descontentamento interno e culpando as instituições multilaterais por ``servirem aos interesses das elites''. Como resultado, o sistema multilateral de comércio está em crise; as dificuldades em combater a evasão fiscal, por exemplo, refletem uma regulamentação internacional inadequada.

A menos que as principais potências econômicas façam esforços sérios para reformar o sistema multilateral, melhorar a cooperação internacional em matéria tributária, facilitar o intercâmbio de tecnologia e enfrentar a crise climática, os países menos desenvolvidos ficarão para trás, deixando suas populações mais vulneráveis arcar com o peso das consequências.

Entre o final da década de 1990 e o final da década de 2000, o apoio público à redistribuição aumentou em quase 70\% das economias avançadas e em desenvolvimento pesquisadas. No entanto, isso não se materializou devido ao maior poder político dos ricos em comparação com os pobres. O recente aumento da desigualdade em todo o mundo está ligado à eliminação de regimes tributários progressivos, à redução do imposto de renda (especialmente para as faixas mais altas) e a outras políticas que beneficiam os segmentos mais ricos da sociedade. Entre as medidas propostas para combater a desigualdade, o relatório identifica a tributação como o principal mecanismo a ser utilizado. Os programas tributários progressivos são um fator importante na redução da desigualdade, enquanto os impostos sobre ganhos de capital e herança também contribuem para reduzir a desigualdade intergeracional.

Em relação às metas tributárias, afirma-se que redistribuir o poder de compra, em vez de renda, pode ser mais eficaz para diminuir a desigualdade de consumo. Uma vez arrecadado por meio de programas tributários, o dinheiro deve ser redistribuído de acordo com políticas públicas. Políticas sociais redistributivas podem incluir benefícios em dinheiro, como subsídios, bem como benefícios não monetários, como programas de saúde e educação pública. Países da América Latina experimentaram com sucesso programas de transferência de renda, e propostas para uma renda básica universal surgiram em diversas partes do mundo. Reformas agrárias também têm um histórico de contribuição para o combate à desigualdade.

Embora ainda não esteja claro qual deve ser a forma dessa mudança estrutural, sabe-se que uma abordagem proativa por parte dos Estados e legisladores é necessária, visto que deixar essas ambições ao mercado não levará à transformação estrutural necessária. Os produtores se concentram no sucesso e no lucro imediatos, muitas vezes ignorando os custos sociais e ambientais. Os consumidores são obrigados a comprar bens e serviços a preços acessíveis. O sistema atual de produção e consumo é impulsionado por mecanismos que não levam em conta a poluição e os custos sociais. As campanhas de conscientização são lentas e inadequadas para a escala e a urgência das questões ambientais. Não existe solução mágica, o que precisa ser discutido são políticas públicas para lidar com as mudanças climáticas e a desigualdade.

O relatório propõe uma mudança no discurso, passando de elogiar a " economia de crescimento mais rápido"{} para elogiar o " país que mais reduziu a desigualdade" . Para que isso seja possível, transformações estruturais não são apenas aceitáveis, mas parte integrante da estratégia para combater a desigualdade. É evidente que as políticas públicas adotadas no passado foram lamentavelmente insuficientes para combater a pobreza e alcançar o desenvolvimento sustentável; o melhor que se pode dizer é que houve alguma estabilidade nos últimos 15 anos, após a deterioração nos 15 anos anteriores.

Entre os impostos sugeridos pelo relatório está um imposto sobre o carbono, com o objetivo de criar um desincentivo financeiro para as emissões de CO2. O relatório também menciona impostos sobre a propriedade e sobre a renda. Durante grande parte do século XX, os países implementaram impostos progressivos (impostos mais altos para indivíduos de renda mais alta) para garantir uma redistribuição mais equitativa e a oferta de serviços públicos de qualidade. Contudo, com base no argumento econômico de que impostos mais baixos incentivam o investimento privado e geram um efeito multiplicador maior do que o investimento público, e também devido ao poder político que indivíduos de alta renda podem exercer, propagando esse argumento, a progressividade tributária diminuiu na maioria das economias desenvolvidas nas últimas quatro décadas, especialmente entre os mais ricos.

Os impostos sobre a propriedade não mitigam diretamente a desigualdade de renda ou riqueza, mas são a principal fonte de receita para financiar a educação nos EUA. Esse tipo de imposto sobre a propriedade é amplamente utilizado em países desenvolvidos e menos em países em desenvolvimento. Os impostos sobre a riqueza, que se aplicam a ativos reais e financeiros, estão se tornando cada vez mais populares. Mas os críticos argumentam que eles geram pouca receita em relação aos custos administrativos e podem ser facilmente contornados por meio de deduções, ocultação de fundos ou realocação de ativos para outros países. Isso levou à eliminação dos impostos sobre a riqueza em muitos contextos. Em 2017, apenas sete países da OCDE mantinham esse tipo de imposto.

O imposto sobre o terreno é considerado o mais eficiente de todos os impostos. No entanto, ainda é uma novidade na maioria dos países. O conceito é simples: uma parcela significativa da riqueza está concentrada no valor da terra urbana, e esse valor aumenta principalmente devido ao investimento público na área. As ações individuais do proprietário de terras não aumentam o valor. O imposto sobre o terreno, portanto, visa o aumento do valor da terra atribuível ao investimento público, um ganho " não merecido"{} para o proprietário.

Como medida de redistribuição de riqueza, existe a proposta de uma renda básica universal. Essa é uma ideia que existe pelo menos desde o final do século XVIII e que ressurgiu no século XXI: um programa que fornece a cada indivíduo, independentemente de quaisquer critérios, dinheiro suficiente para cobrir as necessidades básicas. Os defensores argumentam que isso pode permitir que os indivíduos recusem trabalhos indesejáveis, aumentando o poder de negociação dos trabalhadores e reduzindo a desigualdade.

Os oponentes, no entanto, argumentam que existe um incentivo negativo, funcionando como um desincentivo ao trabalho. É importante notar que a eficácia de qualquer política pública depende de como o programa interage com um sistema de bem-estar social mais amplo. Como complemento, a Renda Básica Universal pode ter um efeito equalizador; no entanto, como substituto de todos os programas, pode agravar a desigualdade. O maior desafio diz respeito aos recursos necessários para financiar o programa. A Finlândia realizou um estudo piloto entre 2017 e 2018, e os resultados preliminares indicam que a Renda Básica Universal não teve efeito sobre os níveis de emprego, mas aumentou a confiança no bem-estar físico e mental dos indivíduos, bem como a confiança nas instituições.

O relatório da ONU já é bastante abrangente. Para um aprofundamento no tema, recomendo o livro de Piketty, " O Capital no Século XXI"\cite{piketty2014capital}. No entanto, como já tenho algumas anotações sobre o relatório do UBS de 2022, vou utilizá-las para fazer alguns comentários adicionais. Este relatório foi elaborado por dois bancos e é a fonte mais atualizada sobre a distribuição de riqueza no mundo; um dos dois bancos se autodeclara o maior gestor de patrimônio global.

Em primeiro lugar devo comentar que grande parte da análise do relatório parece adotar uma metodologia questionável, pois atribui um peso significativo à taxa de câmbio do dólar em relação a outras moedas. Por exemplo, um dos resultados apresentados é que a riqueza global diminuiu, mas grande parte disso\cite{ubs2022wealth} se deve à relação do dólar com outras moedas, conforme relatado. Mesmo assim, há alguns dados que merecem destaque. Primeiro, vemos a confirmação do relatório anterior de que, em grande medida, a redução da desigualdade entre os países é impulsionada pelo crescimento das economias emergentes, especialmente a China, mas também, em grande parte, a Índia. Existem também dados que indicam relativa estabilidade tanto no coeficiente de Gini quanto na riqueza do 1\% mais rico desde os anos 2000.

Em relação à distribuição de riqueza entre diferentes frações da população nos Estados Unidos, a riqueza mediana das famílias negras representa pouco mais de 10\% da riqueza mediana das famílias brancas. A riqueza mediana das famílias hispânicas, em comparação com as famílias brancas, é de 20\%. Como relatado anteriormente, as crises afetam diferentes setores da sociedade de maneiras distintas. Parte da explicação para esse fenômeno reside na natureza diversa da riqueza entre as classes sociais. Por exemplo, a maior parte do grupo hispânico investe em imóveis (ativos físicos ou não financeiros), enquanto os homens brancos, que possuem a maior média e mediana de riqueza, concentram a maior parte de sua riqueza em ativos não físicos (financeiros)\footnote{Os ativos financeiros e não financeiros diferem na forma como são comprados e vendidos. Muitos ativos financeiros, como ações e títulos, são negociados em bolsas de valores e podem ser comprados ou vendidos em qualquer dia útil em que a bolsa esteja aberta. É fácil obter o preço de mercado atual para comprar ou vender esses ativos. Por outro lado, um ativo não financeiro, como um equipamento ou um veículo, pode ser difícil de vender porque não há um mercado ativo de compradores e vendedores. O preço do item não financeiro pode ser incerto, pois não há um padrão de mercado. Em vez disso, muitos ativos não financeiros são vendidos quando o vendedor encontra um comprador em potencial e negocia um preço de venda. O tempo necessário para encontrar um comprador, concluir a venda e transferir o ativo físico torna os ativos não financeiros ilíquidos\cite{chen2026financial,chen2026nonfinancial}.}.

Em números globais, os 52,5\% mais pobres da população detêm apenas 1,2\% de toda a riqueza. Nos países mais pobres, 80\% da população está concentrada na base dessa pirâmide, enquanto nos países mais ricos esse percentual é de 30\%, sendo a distribuição mais transitória em comparação com os países mais pobres. Essa fração da população global também apresenta grande diversidade em todo o mundo. Os 1,1\% mais ricos detêm 45,8\% da riqueza e estão particularmente concentrados em certas regiões e países, compartilhando um estilo de vida e oportunidades semelhantes, com foco em ativos não físicos. Essa distribuição global de riqueza indica a existência de uma pirâmide internacional também entre os países.

Em termos numéricos, é preciso US\$ 9.000 para estar entre os 50\% mais ricos, 16 vezes esse valor para estar entre os 10\% mais ricos (US\$ 140.000) e 125 vezes esse valor para estar entre o 1\% mais rico (US\$ 1.000.000). Retomando a questão da ``redução'' da desigualdade global, as pessoas com riqueza próxima à mediana global tendem a viver em economias emergentes e possuem sua riqueza principalmente na forma de ativos físicos, especificamente imóveis. Durante a pandemia, a população mais pobre foi impactada economicamente, precisando usar suas economias e contrair dívidas. Para a população mais pobre, os ativos não financeiros que possuem representam mais que o dobro da riqueza total, mas o valor desses ativos é compensado pelas dívidas que contraíram. A população mais rica, por outro lado, não só foi menos vulnerável, como também se beneficiou da crise.

O relatório também visa reduzir a questão da desigualdade a duas perguntas:
\begin{itemize}
\item Quão distante está o topo da pirâmide do cidadão médio? Essa pergunta é frequentemente respondida pela porcentagem da riqueza detida pelo 1\% (ou 10\%) mais rico.
\item Quão distante está a base da pirâmide do cidadão médio?
\end{itemize}
Analisando dados de diferentes países, o 1\% mais rico tende a deter entre 30\% e 35\% da riqueza, enquanto os 10\% mais ricos detêm entre 60\% e 65\% da riqueza total. Comparando países, um índice de Gini de 70 é geralmente considerado relativamente baixo e 80, relativamente alto. Lembre-se de que o coeficiente de Gini mede a desigualdade e varia de 0 (igualdade total) a 100 (desigualdade total). A desigualdade aumentou em países com políticas mais liberais e, segundo este relatório, diminuiu na maior parte do mundo devido à crescente importância de ativos não financeiros. No entanto, vale ressaltar que essa redução é tão pequena que eu tenderia a desconsiderá-la.

Voltando aos números globais, os 10\% mais ricos detêm 81\% da riqueza mundial. A distribuição de ativos financeiros e não financeiros não é uniforme. Os ativos não financeiros predominam na base da pirâmide, sendo gradualmente substituídos por ativos financeiros à medida que subimos na pirâmide. O 1\% mais rico concentra sua riqueza principalmente em ativos financeiros. Dessa forma, as variações nos preços dos ativos resultam em mudanças na desigualdade e, consequentemente, na riqueza média e mediana.

Um fato impressionante é que o crescimento da China ao longo de 20 anos (2000 a 2022) é equivalente ao crescimento dos Estados Unidos ao longo de 80 anos (1925-2005), e os próximos cinco anos (2022-2027) devem ser equivalentes a mais 14 anos nos EUA, restando apenas 8 anos de diferença. Nesse ritmo, a China poderá alcançar os EUA até o final da década.

Outro ponto importante é que, durante a pandemia, o setor público socorreu o setor privado por meio do aumento dos gastos governamentais. Entretanto, em 2022 essa mudança temporária na relação entre riqueza pública e privada começou parcialmente a se reverter.

\bibliographystyle{plain}
\bibliography{tese,Artigo1/references,Artigo2/referencias}

\end{document}
